\section{Canonical subsystem marginals under a shared reservoir}
\label{sec:shared-baths}

Two identical systems sharing a reservoir can occupy opposite phases
while each separately has its canonical distribution. Their missing
independence has an exact information-theoretic limit. Thermally coupled
systems with negative specific heat were previously shown to develop
different subsystem phases and to exchange their assignments
in~\cite{RamirezHernandezLarraldeLeyvraz2008}. The results below quantify
the reservoir-size conditions for this behavior in the full microscopic
law and distinguish the accuracy of separate marginals from joint
canonical accuracy.

\subsection{Opposite-phase selection by the exact reservoir}

Let $P_N$ be a one-copy canonical law at the specified inverse temperature
$\beta_N\to\beta>0$. Assume two disjoint measurable phase events
$A_{-,N},A_{+,N}$ partition its state space, with
\begin{equation}
 w_{i,N}:=P_N(A_{i,N})\longrightarrow w_i>0,
 \qquad w_-+w_+=1.
 \label{eq:shared-phase-weights}
\end{equation}
Let $e_{-,N}<e_{+,N}$ be deterministic energy centers, with
$\Delta_N=e_{+,N}-e_{-,N}$, chosen so that
\begin{equation}
 \Delta_N/N\longrightarrow\ell>0,\qquad
 (E_N-e_{i,N})/\sqrt N\mid A_{i,N}\quad\hbox{is tight}.
 \label{eq:shared-tightness}
\end{equation}
Any sequences with these two properties may serve as centers; the
results below use nothing else about them. The assumptions permit
discrete phase laws and a degenerate conditional energy limit. In this
section $N$ denotes the size of one copy.

Take two independent copies under $\Pi_N=P_N\otimes P_N$ and write
$X_1,X_2$ for their complete microscopic states and $S_1,S_2$ for the
phase labels. Their combined energy is $E_N(X_1)+E_N(X_2)$. There is no
direct interaction between the copies. Couple their sum to one exact
reservoir and choose
\begin{equation}
 M_N=e_{-,N}+e_{+,N},\qquad
 \mathcal E_N=M_N+c_N/\beta_N.
 \label{eq:shared-centered-energy}
\end{equation}
The centered weight from~\eqref{eq:centered-weight} is
\begin{equation}
 W_N=
 \begin{cases}
 \exp[\beta_N T_N]
       (1-\beta_N T_N/c_N)^{c_N},&\beta_N T_N<c_N,\\
 0,&\beta_N T_N\ge c_N,
 \end{cases}
 \quad
 T_N=E_N(X_1)+E_N(X_2)-M_N.
 \label{eq:shared-weight}
\end{equation}
Let $z_N=\E_{\Pi_N}W_N$ and $\dd Q_N/\dd\Pi_N=W_N/z_N$.
The bound $0\le W_N\le1$ holds globally, including all interphase and
far-tail energies.

\begin{theorem}[Opposite phases and full marginal accuracy]
\label{thm:shared-selection}
Under~\eqref{eq:shared-phase-weights}--\eqref{eq:shared-tightness},
suppose
\begin{equation}
 N\ll c_N\ll N^2.
 \label{eq:shared-window}
\end{equation}
For $O_N=\{S_1\ne S_2\}$, the centered shared reservoir satisfies
\begin{equation}
 \TV\bigl(Q_N,\Pi_N(\cdot\mid O_N)\bigr)\longrightarrow0.
 \label{eq:shared-selection-tv}
\end{equation}
Consequently, with $Q_N^{(j)}$ denoting either one-copy marginal,
\begin{align}
 \TV(Q_N,\Pi_N)&\longrightarrow w_-^2+w_+^2,
 \label{eq:shared-joint-tv}\\
 \TV(Q_N^{(j)},P_N)&\longrightarrow |w_--1/2|,
 \qquad j=1,2.
 \label{eq:shared-marginal-tv}
\end{align}
\end{theorem}
\begin{proof}
On either opposite-phase assignment $T_N=O_{\Pi_N}(\sqrt N)$
conditionally. The Taylor expansion
$x+\log(1-x)=-x^2/2+O(|x|^3)$, together with $N/c_N\to0$, gives
$\log W_N\to0$ in its conditional probability; the cutoff lies beyond
every fixed fluctuation window. On the same-phase assignments,
$T_N=\pm\Delta_N+O_{\Pi_N}(\sqrt N)$. Because $c_N/N\to\infty$,
the Taylor expansion still applies in conditional probability, but now
\[
 \log W_N=-\frac{\beta_N^2\Delta_N^2}{2c_N}[1+o_{\Pi_N}(1)]
            \longrightarrow-\infty.
\]
Thus $W_N-\ind_{O_N}\to0$ in probability and, by boundedness, in
$L^1(\Pi_N)$. In particular,
$z_N\to p_O:=2w_-w_+>0$. Normalizing the two nonnegative weights
proves~\eqref{eq:shared-selection-tv}.

Conditioning a probability measure on an event of probability $p$
changes it by $1-p$ in total variation. This gives
\eqref{eq:shared-joint-tv}. The exact marginal of
$\Pi_N(\cdot\mid O_N)$ is
\[
 \tfrac12P_N(\cdot\mid A_{-,N})+
 \tfrac12P_N(\cdot\mid A_{+,N}).
\]
The two phase supports are disjoint, so its distance from $P_N$ is
$|w_{-,N}-1/2|$. Total variation contracts under projection, proving
\eqref{eq:shared-marginal-tv}.
\end{proof}

When $w_-=w_+=1/2$, both \emph{complete} subsystem marginals converge
to their canonical laws, while the joint distance tends to $1/2$.
The result does not depend on a density approximation or a bound on the
probability of a particular interfacial configuration. Tightness controls
the atypical canonical mass, and the chosen reservoir cannot amplify
that mass. The second copy accommodates a compensating latent-energy
change, while the ordinary reservoir accommodates the smaller phase
fluctuations.

\subsection{The full microscopic mutual information}
\label{subsec:shared-information}

For $Q\ll P$, define
$\KL(Q\Vert P)=\int\log(\dd Q/\dd P)\dd Q$, with $0\log0=0$.
All logarithms are natural, so information is initially in nats.
For a joint law, its mutual information is
$I_Q(X_1;X_2)=\KL(Q\Vert Q^{(1)}\otimes Q^{(2)})$.
Total variation by itself does not control this quantity on growing
or continuous state spaces. We therefore prove the entropy limit
separately from Theorem~\ref{thm:shared-selection}.

\begin{proposition}[One bit in the full state law]
\label{prop:shared-one-bit}
Under the hypotheses of Theorem~\ref{thm:shared-selection},
\begin{equation}
 I_{Q_N}(X_1;X_2)\longrightarrow\log2.
 \label{eq:shared-one-bit}
\end{equation}
This limit is independent of the original positive phase weights.
\end{proposition}
\begin{proof}
The function $x\log x$ is continuous and bounded on $[0,1]$.
The proof of Theorem~\ref{thm:shared-selection} therefore gives
$\E_{\Pi_N}W_N\log W_N\to0$. Hence
\begin{equation}
 \KL(Q_N\Vert\Pi_N)
 =z_N^{-1}\E_{\Pi_N}W_N\log W_N-\log z_N
 \longrightarrow-\log(2w_-w_+).
 \label{eq:shared-joint-kl}
\end{equation}
Each marginal likelihood $r_{j,N}=\dd Q_N^{(j)}/\dd P_N$ is at most
$1/z_N$, and so is bounded uniformly for all large $N$. By
\eqref{eq:shared-selection-tv} it converges in $L^1(P_N)$ to the
finite-$N$ conditional marginal ratio
\[
 r_N^O=\frac{\ind_{A_{-,N}}}{2w_{-,N}}
             +\frac{\ind_{A_{+,N}}}{2w_{+,N}}.
\]
Both ratios lie in a common compact interval of $[0,\infty)$.
Uniform continuity of $r\log r$ on that interval and convergence
in probability imply convergence of its integrals. Consequently
\begin{equation}
 \KL(Q_N^{(j)}\Vert P_N)
       \longrightarrow\tfrac12\log\frac1{2w_-}
                        +\tfrac12\log\frac1{2w_+}.
 \label{eq:shared-marginal-kl}
\end{equation}
All joint and marginal divergences are finite, because their likelihood
ratios are bounded. The relative-entropy chain rule gives
\begin{equation}
 I_{Q_N}(X_1;X_2)=\KL(Q_N\Vert\Pi_N)
                   -\sum_{j=1}^2\KL(Q_N^{(j)}\Vert P_N).
 \label{eq:shared-kl-chain}
\end{equation}
Substituting~\eqref{eq:shared-joint-kl} and
\eqref{eq:shared-marginal-kl} proves~\eqref{eq:shared-one-bit}.
No microscopic differential entropy is used.
\end{proof}

The limiting phase labels are uniform and perfectly anticorrelated.
Equation~\eqref{eq:shared-one-bit} establishes that no additional
nonzero limiting amount of mutual information remains in their internal
microscopic fluctuations in the window~\eqref{eq:shared-window}.

\subsection{Joint canonical accuracy requires a quadratic-size reservoir}

\begin{corollary}[Optimized joint threshold]
\label{cor:shared-joint-threshold}
Suppose only that, for some centers $a_N$, $(E_N-a_N)/N$ converges
weakly to a law with exactly two distinct positive-weight atoms.
For any fixed number of
independent identical copies $m\ge2$, a shared physical reservoir can
reproduce their complete canonical product law in total variation,
by some total-energy tuning, if and only if $c_N\gg N^2$.
\end{corollary}
\begin{proof}
After a common centering, the total energy divided by $N$ converges
to the $m$-fold convolution of the two-atom law. It has $m+1\ge3$
distinct support points, all of positive weight. Apply
Theorem~\ref{thm:three-support} with $b_N=N$ to the combined subsystem.
The number of copies is fixed, so the choice of $N$ rather than $mN$
does not change the asymptotic condition.
\end{proof}

No phase central limit theorem is required for this corollary. For the
two-copy assumptions above, the centered total energy
$[E_N(X_1)+E_N(X_2)-M_N]/N$ has limiting atoms $-\ell,0,\ell$.
No individual limit for $e_{i,N}/N$ is needed. For the
microscopic models in Section~\ref{sec:microscopic}, an isolated copy
has the sharp threshold $c_N\gg N^{3/2}$, whereas a fixed collection
of at least two copies has the stronger joint threshold $c_N\gg N^2$.
With balanced phase weights, a choice such as $c_N=N^{5/4}$ nevertheless
gives canonical one-copy marginals under the shared bath. The isolated
copy and the shared pair are different composite systems; the phase
anticorrelation records the role of the additional energy store.
Two extensions are given in Appendix~\ref{app:shared-extensions}:
selection of a fixed number of plus-phase copies among $m$ copies, and
the correlated within-phase fluctuations that appear when the capacity
is only proportional to $N$.

\subsection{Balancing the reference phase weights}
\label{subsec:shared-balancing}

The limiting phase weights at a fixed coexistence temperature need not
be equal. To realize the canonical-marginal case in a microscopic model,
the canonical reference itself can be chosen at a specified nearby
temperature. The following elementary statement makes the change of
reference explicit.

\begin{lemma}[A bounded-energy temperature shift]
\label{lem:shared-temperature-shift}
Let the spin energy $U_N/N$ be uniformly bounded and suppose the phase
conditional laws at $\beta_c$ concentrate at $u_-<u_+$, with positive
limiting weights $w_-,w_+$. For
$\beta_N=\beta_c+\delta/N+o(N^{-1})$, the new limiting weight ratio is
\begin{equation}
 \frac{w_-^{(\delta)}}{w_+^{(\delta)}}
       =\frac{w_-}{w_+}\exp[\delta(u_+-u_-)].
 \label{eq:shared-temperature-ratio}
\end{equation}
The new spin law conditional on either phase approaches the old
phase-conditional law in total variation. Hence any conditional weak
fluctuation limit, with the same center and scale, persists.
\end{lemma}
\begin{proof}
The change of canonical measure has an unnormalized likelihood
$\exp[-(\beta_N-\beta_c)U_N]$. It is bounded above and below by
positive constants uniformly in $N$. On phase $i$ it converges in
probability to $e^{-\delta u_i}$. Bounded convergence gives the
limiting phase normalizations and~\eqref{eq:shared-temperature-ratio}.
After normalizing within either phase, the likelihood converges to
one in $L^1$, proving the conditional total-variation assertion.
\end{proof}

Writing $\ell=u_+-u_-$, asymptotic balance is obtained by
\begin{equation}
 \beta_N=\beta_c-\frac{\log(w_-/w_+)}{N\ell}+o(N^{-1}).
 \label{eq:shared-balanced-temperature}
\end{equation}
The negative sign is consistent with warming a system whose ordered,
lower-energy phase initially has excess weight. For the short-range
model, $w_-/w_+=q$, and the required shift is
$-\log q/(N\ell)$. For the three-state mean-field model, the ratio is
$3r_*$ from~\eqref{eq:mf-phase-weights} and $\ell=J/12$.
Both spin Hamiltonians are bounded per particle, so
Lemma~\ref{lem:shared-temperature-shift} applies directly. For the
short-range positive contour decomposition, retain the conditional bond
kernel $K_{\beta_c}(\dd A\mid\sigma)$ from the Edwards--Sokal law at
$\beta_c$ and define the auxiliary law
\[
 \widehat P_N(\dd\sigma,\dd A)
   =P_{\beta_N}(\dd\sigma)K_{\beta_c}(\dd A\mid\sigma).
\]
Its likelihood relative to the old joint law is precisely
$\dd P_{\beta_N}/\dd P_{\beta_c}$, bounded above and below by positive
constants. Thus the fixed contour events still give a positive
decomposition of the new physical spin law. Their conditional
exponential moments remain bounded, and their exceptional mass increases
by at most a constant factor. The retained bond kernel is auxiliary;
it need not be the Edwards--Sokal kernel at $\beta_N$. The same bounded
spin likelihood directly preserves the mean-field positive phase moment
bounds. Hence the isolated-copy $N^{3/2}$ thresholds persist under these
shifts.

If independent momenta are included, their canonical partition factor
is common to both spin phases and cancels from the weight ratio. Use
the kinetic law at the same $\beta_N$ as the spin reference. To check
the effect on the complete kinetic law, its relative entropy against
the law at $\beta_c$ is
\[
 aN\left[\frac{\beta_c}{\beta_N}-1
                     -\log\frac{\beta_c}{\beta_N}\right]=O(N^{-1}).
\]
This follows either from the product Gaussian momenta or their Gamma
energy law. Pinsker's inequality gives total-variation convergence of
the two kinetic laws. The mean energy changes by $O(1)$ and the
$\sqrt N$ Gaussian fluctuation limit is unchanged. No globally bounded
likelihood is asserted for the unbounded kinetic energy.

Use the resulting balanced reference for every copy and the same
$\beta_N$ in~\eqref{eq:shared-centered-energy}. Then
Theorem~\ref{thm:shared-selection} gives two canonical subsystem
marginals, joint total-variation error $1/2$, and
Proposition~\ref{prop:shared-one-bit} gives one bit of mutual
information. The reference being reproduced is the balanced canonical
law at $\beta_N$, not the unbalanced law at $\beta_c$.

\subsection{Phase events for the microscopic models}
\label{subsec:shared-microscopic}

The hypotheses \eqref{eq:shared-phase-weights}--\eqref{eq:shared-tightness}
ask for a partition of the state space into two phase events with
conditionally tight $\sqrt N$ energy fluctuations. For the mean-field
model of Section~\ref{subsec:meanfield}, take the ordered and
disordered Voronoi cells; Lemma~\ref{lem:mf-exponential} gives
conditional exponential moments on the $\sqrt N$ scale, hence
tightness, with or without momenta. For the short-range model of
Section~\ref{subsec:short-range}, take the midpoint energy events
$B_{o,N}$ and $B_{d,N}$ defined in
Appendix~\ref{app:short-range-proof}; their conditional standardized
energies converge weakly by \eqref{eq:sr-conditional-clt}, hence are
tight. The contour decomposition with its exceptional mass is not
needed for the shared-bath theorems, because the centered weight
\eqref{eq:shared-weight} is bounded by one and cannot amplify rare
configurations. After the balancing shift of
Section~\ref{subsec:shared-balancing}, Lemma~\ref{lem:shared-temperature-shift}
preserves both partitions and their conditional limits.
